% ---------------------------------------------------------------------------
%  Preambolo comune. Tenuto separato da main.tex cosi' che l'indice dei
%  capitoli resti leggibile in un colpo d'occhio.
% ---------------------------------------------------------------------------

% --- lingua e codifica -----------------------------------------------------
\usepackage[T1]{fontenc}
\usepackage[utf8]{inputenc}
\usepackage[english]{babel}
\usepackage{lmodern}
\usepackage{microtype}

% --- impaginazione ---------------------------------------------------------
\usepackage[a4paper,inner=3cm,outer=2.5cm,top=2.8cm,bottom=2.8cm]{geometry}
\usepackage{fancyhdr}
\usepackage{emptypage}
\setlength{\parskip}{0.4em}

\setlength{\headheight}{14pt}   % richiesto da fancyhdr
\pagestyle{fancy}
\fancyhf{}
\fancyhead[LE,RO]{\thepage}
\fancyhead[RE]{\nouppercase{\leftmark}}
\fancyhead[LO]{\nouppercase{\rightmark}}
\renewcommand{\headrulewidth}{0.4pt}
\fancypagestyle{plain}{\fancyhf{}\fancyfoot[C]{\thepage}%
  \renewcommand{\headrulewidth}{0pt}}

% --- matematica ------------------------------------------------------------
\usepackage{amsmath}
\usepackage{amssymb}

% --- unita' di misura ------------------------------------------------------
% Un solo punto di configurazione per tutte le unita' del report: cosi' non
% si finisce con "30 A_RMS", "30 Arms" e "30 A rms" in tre capitoli diversi.
%
% NOTA SULLA SCELTA. Lo strumento naturale sarebbe siunitx, ed e' quello per
% cui la sintassi qui sotto e' scritta. Su questa macchina non e' usabile:
% MiKTeX ha siunitx 3.5.12 (2026-09-08) su un kernel LaTeX2e 2023-11-01 con
% L3 2024-01-04, e il pacchetto usa variabili che quel kernel non definisce.
% Fallisce anche su un documento minimo di quattro righe, quindi non e' un
% problema di questo report.
%
% Per non lasciare il documento ostaggio di un toolchain rotto, le macro
% sono reimplementate qui in LaTeX puro. I punti di chiamata restano
% identici a siunitx: quando MiKTeX sara' aggiornato bastera' cancellare
% questo blocco e togliere il commento alle due righe in fondo.
\usepackage{textcomp}

\makeatletter
% I simboli sono definiti GLOBALMENTE, non solo dentro \si. Serve perche'
% \caption scrive i propri argomenti negli elenchi di figure e tabelle cosi'
% come sono: la' dentro \volt viene riletto a livello superiore, e se fosse
% definito solo localmente sarebbe una sequenza non definita.
% prefissi
\def\yocto{y}\def\pico{p}\def\nano{n}\def\micro{\textmu}\def\milli{m}
\def\centi{c}\def\deci{d}\def\kilo{k}\def\mega{M}\def\giga{G}
% unita' di base e derivate usate nel report
\def\metre{m}\def\meter{m}\def\second{s}\def\gram{g}\def\kelvin{K}\def\minute{min}
\def\ampere{A}\def\volt{V}\def\watt{W}\def\joule{J}\def\hertz{Hz}
% L'ohm NON usa \textohm: il glifo TS1 ha una mappatura ToUnicode rotta e
% l'estrazione di testo dal PDF restituisce "W" al posto di "ohm", rendendo
% inutilizzabili ricerca e copia-incolla ("40 mohm" diventava "40 mW").
% L'Omega matematico e' dritto per default in TeX ed e' mappato bene.
\def\ohm{\ensuremath{\Omega}}\def\farad{F}\def\henry{H}\def\coulomb{C}
\def\siemens{S}\def\celsius{\textcelsius}\def\percent{\%}
\def\VA{VA}\def\var{var}
\def\degree{\textdegree}\def\rpm{rpm}
% raffreddamento (11/10/2026): portate, pressioni, volumi. Litro maiuscolo:
% la l minuscola accanto a un numero si legge 1 ("10 l/min" sembrava "101/min")
\def\litre{L}\def\liter{L}\def\pascal{Pa}\def\hour{h}\def\decibel{dB}\def\cubic#1{#1\ensuremath{^{3}}}
% grandezze efficaci e di picco: scritte sempre allo stesso modo
\def\Arms{A\textsubscript{rms}}\def\Vrms{V\textsubscript{rms}}
\def\Apk{A\textsubscript{pk}}\def\Vpk{V\textsubscript{pk}}
% separatore
\def\per{/}
\def\squared{\ensuremath{^{2}}}\def\cubed{\ensuremath{^{3}}}
\let\pecd@units\relax

% Tutte protette: compaiono dentro \caption, che scrive negli elenchi di
% figure e tabelle. Senza \protected verrebbero espanse durante quella
% scrittura, dove i simboli di \pecd@units non sono definiti.
% Numero: \ensuremath cosi' che \pm, \times e simili funzionino.
\protected\def\num#1{\ensuremath{#1}}
% Unita' sola.
\protected\def\si#1{{\pecd@units\textup{#1}}}
% Parte unita' di \SI: se vuota non mette lo spazio fine, cosi' si possono
% scrivere intervalli come \SI{10}{}--\SI{30}{\kilo\hertz}.
\protected\def\pecd@unit#1{%
  \edef\pecd@tmp{\detokenize{#1}}%
  \ifx\pecd@tmp\@empty\else\,\si{#1}\fi}
\protected\def\SI#1#2{\num{#1}\pecd@unit{#2}}
\protected\def\qty#1#2{\SI{#1}{#2}}
\protected\def\numrange#1#2{\num{#1}\,--\,\num{#2}}
\protected\def\SIrange#1#2#3{\num{#1}\,--\,\SI{#2}{#3}}
% Accettate e ignorate, per compatibilita' con la sintassi siunitx.
\newcommand{\sisetup}[1]{}
\newcommand{\DeclareSIUnit}[2]{}
\makeatother

% Quando MiKTeX sara' aggiornato, sostituire tutto il blocco qui sopra con:
%   \usepackage{siunitx}
%   \sisetup{detect-all, per-mode=symbol}
%   \DeclareSIUnit{\Arms}{A_{\mathrm{rms}}}  ... e le altre tre

% --- figure e tabelle ------------------------------------------------------
\usepackage{graphicx}
\graphicspath{{figures/}}
\usepackage{booktabs}
\usepackage{tabularx}
\usepackage{multirow}
\usepackage{longtable}
\usepackage[font=small,labelfont=bf]{caption}
\usepackage{subcaption}
\usepackage{float}
\usepackage{rotating}   % sidewaysfigure: fogli dello schema in orizzontale (cap. 9)

% Pagine di soli float: in alto e a distanza fissa. Di default LaTeX le
% centra in verticale e lascia mezzi fogli bianchi fra una tabella e una figura.
\makeatletter
\setlength{\@fptop}{0pt}
\setlength{\@fpsep}{2\baselineskip}
\setlength{\@fpbot}{0pt plus 1fil}
\makeatother

% colonna centrata a larghezza automatica per le tabelle numeriche
% Colonne di testo allineate a sinistra: una cella su piu' righe centrata si
% legge male (revisione dell'11/10/2026).
\newcolumntype{Y}{>{\raggedright\arraybackslash}X}
\newcolumntype{L}[1]{>{\raggedright\arraybackslash}p{#1}}

% --- colori e riquadri -----------------------------------------------------
\usepackage[table]{xcolor}
\definecolor{notebg}{RGB}{245,245,245}
\definecolor{noteline}{RGB}{120,120,120}
\usepackage[most]{tcolorbox}

% Riquadro per le discrepanze trovate durante il lavoro. Sono il valore
% aggiunto del progetto rispetto a un'esecuzione a ricetta, quindi hanno un
% formato riconoscibile invece di sparire nel corpo del testo.
\newtcolorbox{finding}[1]{%
  colback=notebg, colframe=noteline, boxrule=0.4pt, arc=1pt,
  left=6pt, right=6pt, top=5pt, bottom=5pt,
  title={\textbf{#1}}, fonttitle=\small, coltitle=black,
  colbacktitle=notebg, breakable}

% --- codice ----------------------------------------------------------------
% Schemi disegnati (figures/tikz/): TikZ e' gia' caricato da tcolorbox.
% Palette uguale a quella delle figure Python (config.COLORI): leggibile in
% scala di grigi e dai daltonici.
\usepackage[european,straightvoltages]{circuitikz}
\usepackage{pgfplots}
\pgfplotsset{compat=1.17}
\usetikzlibrary{arrows.meta,positioning,calc,fit,backgrounds,decorations.pathreplacing,shapes.geometric,shapes.misc}
\definecolor{fblu}{HTML}{0072B2}
\definecolor{farancio}{HTML}{D55E00}
\definecolor{fverde}{HTML}{009E73}
\definecolor{fgrigio}{HTML}{777777}

% --- tabelle leggibili (revisione del 18/10/2026) --------------------------
% \zebra prima di un tabular: righe alterne grigio chiaro (colortbl, possibile
% dopo l'aggiornamento di MiKTeX del 18/10/2026), un po' piu' di aria fra le
% righe. Le celle PASS/FAIL hanno lo sfondo dell'esito.
\newcommand{\zebra}{\rowcolors{2}{white}{black!6}\renewcommand{\arraystretch}{1.15}}
\newcommand{\pass}[1]{\cellcolor{fverde!30}#1}
\newcommand{\fail}[1]{\cellcolor{farancio!30}#1}
\newcommand{\aperto}[1]{\cellcolor{yellow!45}#1}
\newcommand{\nd}{\cellcolor{black!10}--}
\usepackage{listings}
\lstset{
  basicstyle=\ttfamily\footnotesize,
  breaklines=true,
  frame=single,
  rulecolor=\color{noteline},
  columns=fullflexible,
  keepspaces=true
}

% --- bibliografia ----------------------------------------------------------
\usepackage{csquotes}   % richiesto da biblatex quando c'e' babel
\usepackage[backend=biber,style=ieee,sorting=none]{biblatex}
\addbibresource{references.bib}

% --- riferimenti incrociati ------------------------------------------------
\usepackage[hidelinks]{hyperref}
\hypersetup{
  pdftitle  = {Three-Phase Two-Level SiC Traction Inverter},
  pdfauthor = {Matteo Borghi},
  pdfsubject= {Automotive Power Electronics Circuit Design, Project 7},
  bookmarksnumbered = true
}
\usepackage[capitalise,noabbrev]{cleveref}
\crefname{chapter}{Chapter}{Chapters}
\crefname{section}{Section}{Sections}
\crefname{figure}{Figure}{Figures}
\crefname{table}{Table}{Tables}
\crefname{equation}{Equation}{Equations}
% dopo \appendix i capitoli sono appendici: "Appendix A", non "Chapter A"
\crefname{appendix}{Appendix}{Appendices}

% --- macro di progetto -----------------------------------------------------
% I nomi ricorrenti si scrivono una volta sola: se cambia il dispositivo,
% cambia qui e non in trenta punti del testo.
\newcommand{\device}{E3M0016120K}
\newcommand{\driver}{UCC21750}
\newcommand{\isoamp}{AMC1302}
\newcommand{\isoampv}{AMC1311}
\newcommand{\dcdc}{MGJ2D121505SC}
\newcommand{\Rthjc}{R_{\theta,\mathrm{JC}}}
\newcommand{\Rthcs}{R_{\theta,\mathrm{CS}}}
\newcommand{\Rthha}{R_{\theta,\mathrm{HA}}}
\newcommand{\Tj}{T_{\mathrm{j}}}
\newcommand{\Tcase}{T_{\mathrm{case}}}
\newcommand{\Ths}{T_{\mathrm{hs}}}
\newcommand{\Rdson}{R_{\mathrm{DS(on)}}}
\newcommand{\Eon}{E_{\mathrm{on}}}
\newcommand{\Eoff}{E_{\mathrm{off}}}
\newcommand{\fsw}{f_{\mathrm{sw}}}
\newcommand{\Vdc}{V_{\mathrm{DC}}}

% Segnaposto visibile per i capitoli non ancora scrivibili.
\newcommand{\stubnote}[1]{%
  \begin{center}
  \fcolorbox{noteline}{notebg}{%
    \parbox{0.9\textwidth}{\small\textbf{Placeholder.} #1}}
  \end{center}}
